\subsection{Cinemática Diferencial e Jacobiano}

A cinemática diferencial descreve como as velocidades de junta geram as velocidades da ferramenta em um referencial escolhido.
Foi utilizada a matriz Jacobiana $\mathbf{J}(\vec q)$ para mapear
$\dot{\vec q}$ em velocidade linear e angular da ferramenta.
% Este mapeamento é a base para:
% (i) controle no espaço de tarefa,
% (ii) análise de singularidades e
% (iii) conversão entre forças cartesianas e torques de junta.

\subsubsection{Velocidades cartesianas e Jacobiana}
A partir do referencial da base $\{0\}$, a velocidade (linear e angular) cartesiana da ferramenta é obtida diretamente a partir das velocidades das juntas pela Jacobiana geométrica.
Para o MR configuração 4R1P, no mesmo sistema de coordenadas, as contribuições de cada junta ao movimento da ferramenta:
\begin{equation}
\label{eq:xdot_J_qdot}
{}^{0}\!\dot{\vec x}_{E} \;=\;
\begin{bmatrix}
{}^{0}\!\vec v_{E} \\
{}^{0}\!\vec \omega_{E}
\end{bmatrix}
= \mathbf{J}(\vec q)\,\dot{\vec q},
\end{equation}
em que $\dot{\vec q}$ reúne as velocidades.

\subsubsection{Construção do Jacobiano}
As colunas do Jacobiano são construídas no referencial $\{0\}$ a partir dos eixos
${}^{0}\!\vec z_{\,i-1}$ e
dos vetores posição ${}^{0}\!\vec p_{\,\cdot}$ obtidos pela cinemática direta.

Para a junta revoluta:
\begin{equation}
\begin{cases}
J_{v,i} =
{}^{0}\!\vec{z}_{\,i-1} \times 
\left({}^{0}\!\vec{p}_{\,E} - {}^{0}\!\vec{p}_{\,i-1}\right), \\
J_{\omega,i} =
{}^{0}\!\vec{z}_{\,i-1}.
\end{cases}
\end{equation}

Para a junta prismática:
\begin{equation}
\begin{cases}
J_{v,i} = {}^{0}\!\vec{z}_{\,i-1}, \\
J_{\omega,i} = \vec 0.
\end{cases}
\end{equation}

Portanto, cada coluna de $\mathbf{J}$ representa a contribuição de cada junta para o movimento linear e angular do MR.

\subsubsection{Efeito da Base}
Quando a base está fixa, usa-se diretamente
\eqref{eq:xdot_J_qdot}.
Para base flutuante, a velocidade da ferramenta é a soma do efeito do movimento da base com o efeito do movimento das juntas:
\begin{equation}
{}^{0}\!\dot{\vec x}_{E} =
\mathbf{J}_{base}(\cdot)\,{}^{0}\!\dot{\vec x}_{B}
\;+\;
\mathbf{J}(\vec q)\,\dot{\vec q}.
\end{equation}
Neste trabalho, $\mathbf{J}_{base}(\cdot)$ é tratada de forma agregada na seção de acoplamento base--manipulador; aqui, assume-se base fixa no uso de $\mathbf{J}(\vec q)$ para controle.

\subsubsection{Mapeamento de torques e forças}
O mapeamento entre as forças--momentos (\emph{wrench}) aplicados na ferramenta e os torques nas juntas é dado por:

\begin{equation}
\vec\tau = \mathbf{J}^T (\vec q)\,\vec w_{E},
\end{equation}

em que a \emph{wrench} na ferramenta é
\begin{equation}
\vec w_{E} =
\begin{bmatrix}
\vec f_{E} \\
\vec m_{E}
\end{bmatrix},
\end{equation}
com $\vec f_{E}$ (N) e $\vec m_{E}$ (N$\cdot$m) expressos no mesmo referencial e aplicados no mesmo ponto $E$.
Portanto, ($\mathbf{J}^{\top}\vec w_{E}$) transforma as forças--momentos cartesianos desejados (ou medidos) na ferramenta em torques equivalentes nas juntas.

\subsubsection{Singularidades}
Uma singularidade ocorre quando pequenas variações desejadas no efetuador $\{E\}$ exigem velocidades de junta excessivas em alguma direção.
Isso acontece quando algum termo de menor valor de $\mathbf{J}(\vec q)$ tende a zero e o número de condição cresce.
A pseudo-inversa amortecida pode ser utilizada, e é expressa por:

\begin{equation}
\label{eq:pseudoinversa}
\dot{\vec q} =
\mathbf{J}^T\,\big(\mathbf{J}\,\mathbf{J}^T
+ \lambda^2 I\big)^{-1}\,{}^{0}\!\dot{\vec x}_{E,d},
\end{equation}
com $\lambda>0$.
% Para o MR deste trabalho, de configuração 4R1P, a dimensão de tarefa é $5$ (posição 3D e dois ângulos relevantes à captura), o que reduz o risco de singularidades.